\section{The Khovanov arc}\label{sec:khovanov}

Throughout this section $q$ is odd, $T_q=Q_{1/3}+Q_{1/q}$, and for $r\in\Q\cup\{\infty\}$
\[
 K_r(q)=\operatorname{num}\bigl(Q_r+Q_{1/3}+Q_{1/q}\bigr),
\]
so that $K_{-1/2}(q)=P(-2,3,q)$ and $K_{-3/4}(7)=K_*$.  For twisted complexes $X,Y$ over $\mathcal B$
we write $\HF(X,Y)$ for the homology of $\operatorname{Mor}(X,Y)$, as in
Lemma~\ref{lem:wrapped-reduction}; for curves in pairing position it is their Lagrangian Floer
homology \cite[Thm.~5.22]{KWZ}.  The object $(N,b_{S_{18}})$ of Corollary~\ref{cor:aux-conditional}
is the sum of a rational arc and a closed curve carrying a two-dimensional local system, and one
further Maurer--Cartan term, supported near the corner $(0,0)$, carries it to the Bar-Natan invariant
of $T_q$ (Theorem~\ref{thm:structure-full}).  No rational closure distinguishes the two objects
(Corollary~\ref{cor:closures-full}).

\subsection{Khovanov invariants in Smith's coordinates}\label{ss:kh-smith}
Let $G$ be the group of isometries of $\R^2$ generated by the translations by $2\pi\Z^2$ and the
reflection $v\mapsto-v$, so that the pillowcase is $\R^2/G$ in the coordinates $(\gamma,\theta)$ of
\S\ref{ss:models}.  The group $G$ acts freely on the complement of the \emph{lattice} $\pi\Z^2$, and
the quotient map $\R^2\setminus\pi\Z^2\to\Pst$ is a covering, which we call the \emph{planar cover}.
The lattice point $(m\pi,n\pi)$ lies over the corner whose coordinates are the residues of $m\pi$
and $n\pi$ modulo $2\pi$.  A \emph{lattice line} is a line in $\R^2$ through two lattice points.  The
lattice lines parallel to a given one form two $G$-orbits: the lines of one orbit pass over two of
the corners, and those of the other over the remaining two.  Every straight arc between two corners,
such as the line of a rational tangle (\S\ref{ss:models}), is the image of a segment joining
consecutive lattice points of a lattice line.

The translation $\Phi(\gamma,\theta)=(\gamma+\pi,\theta)$ normalizes $G$ and therefore descends to an
involution of the pillowcase.  It is a symplectomorphism of $\Pst$, it exchanges the corners
$(0,0)\leftrightarrow(\pi,0)$ and $(\pi,\pi)\leftrightarrow(0,\pi)$, and it is the translation of
Remark~\ref{rem:sign}.

Kotelskiy, Watson and Zibrowius assign to a pointed four-ended tangle $T$ two immersed multicurves
with local systems in the four-punctured sphere: the Bar-Natan invariant $\widetilde{\mathrm{BN}}(T)$,
and the Khovanov invariant $\widetilde{\Kh}(T)$, which is the mapping cone of $H\cdot\mathrm{id}$ on
$\widetilde{\mathrm{BN}}(T)$ \cite[Def.~6.5]{KWZ}.  Both are objects of $\operatorname{Tw}(\mathcal B)$,
and their arc components end at the three non-special punctures.  Under the identification of
\cite[\S8.2]{KWZ} the special puncture is the corner $(0,\pi)$, which is the corner that
holonomy-perturbed tangle character varieties miss in the coordinates of CHKK
\cite[Prop.~11.1 and \S11.6]{CHKK}.  In Smith's coordinates that corner is $(\pi,\pi)$
(\S\ref{ss:models}).  We therefore draw an invariant $Y$ of \cite{KWZ} in Smith's coordinates as
$\Phi(Y)$, and we set
\[
 \mathrm{BN}_q:=\Phi\bigl(\widetilde{\mathrm{BN}}(T_q)\bigr).
\]
For a rational tangle $Q$, the arc $\widetilde{\mathrm{BN}}(Q)$ and the figure-eight curve
$\widetilde{\Kh}(Q)$ around it agree, under this identification, with the character varieties
$\Rp_\pi(Q)$ and $\Rnat_\pi(Q)$ \cite[\S1.8 and \S8.2]{KWZ}.  When $Q_r$ is glued to $T_q$, its
mirror $Q_{-r}$ enters the pairing theorem \cite[Thm.~1.9]{KWZ}, and in Smith's coordinates the curve
$\Phi\bigl(\widetilde{\Kh}(Q_{-r})\bigr)$ is the earring curve $E_r$ of $\Qh{r}$ used in
Section~\ref{sec:smith}, a Herald--Kirk figure-eight \cite[Def.~6.7]{HK}; we have verified this
identification directly, by comparing encodings, for 49 slopes.  We call the straight arc that $E_r$
surrounds the \emph{arc of $E_r$}.  For $r=m/n$ in lowest terms it joins $(\pi,0)$ to $(0,\pi)$ if $m$
and $n$ are odd, $(0,0)$ to $(0,\pi)$ if $n$ is even, and $(0,0)$ to $(\pi,0)$ if $m$ is even
(\S\ref{ss:models}).  The complexes $\mathcal E$ and $\mathcal E_*$ of
Section~\ref{sec:smith} encode $E_{-1/2}$ and $E_{-3/4}$.  Since $\Phi$ is a symplectomorphism, the
pairing theorem therefore reads
\begin{equation}\label{eq:kwz-pairing}
 \dim_{\F_2}\HF\bigl(E_r,\mathrm{BN}_q\bigr)=\rk\widetilde{\Kh}\bigl(K_r(q);\F_2\bigr),
 \qquad r\in\Q\cup\{\infty\}.
\end{equation}
We have compared the two sides of \eqref{eq:kwz-pairing} at the 340 slopes used for $q=5$ and the 342
slopes used for $q=7$ in Corollary~\ref{cor:closures-full}, computing the right side with Khoca
\cite{Khoca}; they agree at every slope.

\subsection{The structure theorem}\label{ss:kh-structure}
Put $r_q=(1-q)/(3q-4)$.  Its numerator is even and its denominator is odd, so the arc of $E_{r_q}$
joins $(0,0)$ to $(\pi,0)$.  We use three curves.
\begin{itemize}
\item $\alpha_q$ is the arc of $E_{r_q}$, an embedded arc from $(0,0)$ to $(\pi,0)$.
\item $c$ is the boundary of a regular neighborhood of the arc of $E_{-1}$, which joins $(\pi,0)$ to
 $(0,\pi)$.  Thus $c$ is a simple closed curve separating $\{(\pi,0),(0,\pi)\}$ from
 $\{(0,0),(\pi,\pi)\}$, and its lifts are lines parallel to the lifts of the arc of $E_{-1}$ that
 pass through no lattice point.
\item $(c,J_2)$ is $c$ with the indecomposable two-dimensional local system
 $J_2=\left(\begin{smallmatrix}1&1\\0&1\end{smallmatrix}\right)$.  As a twisted complex it is the
 doubled word of $c$: the doubled word represents $c$ with the transposition local system, which over
 $\F_2$ is similar to $J_2$ \cite[\S5.3]{KWZ}, \cite[Lem.~2.24]{AurouxSmith}.
\end{itemize}
Up to isotopy, $\Phi$ carries each of these curves to itself.  The arc $\alpha_q$ and the arc
$\mathrm{BN}_q$ are bigraded, as are all invariants of \cite{KWZ}.  The curve $c$ admits a
$\delta$-grading in the sense of \cite[Def.~4.44]{Zibrowius} but no bigrading, and neither does
$(c,J_2)$.  Indecomposable two-dimensional local systems also occur on components of
$\widetilde{\mathrm{BN}}$ of other tangles \cite[Rem.~3.20]{KWZmut}.

Write $N_q$ for the reduced twisted complex that encodes $L_{q,\mathrm{PL}}$ in the chart of
\S\ref{ss:explicit}, so that $N_7=N$; the complex $N_5$ is displayed in \eqref{eq:N5}.  Each is an arc
from $(0,0)$ to $(\pi,0)$, and its end at $(0,0)$ is the generator $36$ of $N_7$, respectively $0$ of
$N_5$.  In that chart the faces $R$, $M$, $L$ contain the corners $(0,0)$, $(\pi,0)$, $(\pi,\pi)$, so
an arrow labelled $R$ is a chord around the corner $(0,0)$.  Two reduced twisted complexes are
\emph{strictly isomorphic} if an idempotent-preserving bijection of their generators carries one
differential to the other entry by entry.  The four-arrow element $b_5\in\operatorname{End}(N_5)$ of
\eqref{eq:b5} satisfies $D_{N_5}(b_5)=b_5^2=0$, and $(N_5,b_5)$ is strictly isomorphic to the encoding
of the census smoothing of $L_{5,\mathrm{PL}}$ at each of two of its self-intersection orbits; it
raises the pairing with $E_{-1/2}$ from $5$ to $7=q+2$.

To describe $\mathrm{BN}_q$, fix a lift $p_0$ of $(0,0)$ and let $p_M$ be the lift of $(\pi,0)$ for
which the segment $[p_M,p_0]$ is a lift of $\alpha_q$.  Let $\ell$ be the lattice line through $p_0$
parallel to the lifts of $c$; its lattice points lie alternately over $(0,0)$ and $(\pi,\pi)$.  Number them $p_j$,
$j\in\Z$, consecutively and so that the angle between $p_1-p_0$ and $p_0-p_M$ is acute, and let
$\ell^+$ be the open half-plane bounded by $\ell$ that does not contain $p_M$.  The translation by
$p_2-p_0\in2\pi\Z^2$ generates the stabilizer in $G$ of each lift of $c$.

\begin{theorem}[Theorem~\ref{thm:structure}]\label{thm:structure-full}
\begin{enumerate}
\item[(a)] Let $q=7$ and $b\in\{b_{S_{18}},b_{S_{25}}\}$, or $q=5$ and $b=b_5$.  Then $(N_q,b)$ is
 strictly isomorphic to $\alpha_q\oplus(c,J_2)$.
\item[(b)] Let $5\le q\le21$ be odd.  Then $\mathrm{BN}_q$ is the arc whose lift starts at $p_M$, runs
 along $[p_M,p_0]$, crosses $\ell$ between $p_{-1}$ and $p_0$, and then runs in $\ell^+$, parallel to
 $\ell$ and close to it, until it ends at $p_4$.
\item[(c)] Let $q$ and $b$ be as in (a), and let
 \[
  \kappa_7=0\xrightarrow R19+36\xrightarrow R19\in\operatorname{End}(N_7),\qquad
  \kappa_5=12\xrightarrow R11+0\xrightarrow R11\in\operatorname{End}(N_5).
 \]
 Then $(\delta_{N_q}+\kappa_q)^2=(\delta_{N_q}+b+\kappa_q)^2=0$, and $(N_q,b+\kappa_q)$ is strictly
 isomorphic to $\mathrm{BN}_q$.
\end{enumerate}
\end{theorem}

Part (b) says that $\mathrm{BN}_q$ arises from $\alpha_q$ by sliding its end at $(0,0)$ twice around
that corner, parallel to $c$: beyond $p_0$ the lift covers two periods of a lift of $c$, ending at
the lift $p_4=p_0+2(p_2-p_0)$ of $(0,0)$.  In (c), the arrows $0\xrightarrow R19$ and
$12\xrightarrow R11$ are terms of $\delta_{N_7}$ and $\delta_{N_5}$, so adding $\kappa_q$ deletes one
arrow and adds an arrow labelled $R$ at the end of $N_q$ at $(0,0)$.  Both terms of $\kappa_q$ are
chords around the corner $(0,0)$.

\begin{proof}
Each statement asserts that two explicit reduced twisted complexes are strictly isomorphic, or that
an explicit matrix over $\mathcal B$ squares to zero.  The complexes
of $\alpha_q$, of $c$, of its doubled word and of the arc in (b) are obtained by encoding their lifts
in the chart of \S\ref{ss:explicit}; $\widetilde{\mathrm{BN}}(T_q)$ is computed with Zibrowius's
program \cite{Zibkht} and drawn in Smith's coordinates by $\Phi$.  All complexes involved are of loop
type, that is, every generator meets at most two arrows, so a strict isomorphism is decided by
comparing the cyclic or linear words of the components, matching idempotents, labels and the
directions of arrows.

For (a), $(N_7,b_{S_{18}})$ and $(N_7,b_{S_{25}})$ have components of $23$ and $8$ generators, and
$(N_5,b_5)$ of $15$ and $8$; in each case the first component is strictly isomorphic to $\alpha_q$
and the second to the doubled word of $c$, which represents $(c,J_2)$.  For (b), the encoding of the
arc described there is strictly isomorphic to $\mathrm{BN}_q$ for $q=5,7,\dots,21$; it is the only
one among the $24$ arcs obtained from $\alpha_q$ by sliding one of its ends along $c$, in either
direction and on either side, through one, two or three periods.  For (c), direct multiplication in
$\mathcal B$ gives the two Maurer--Cartan identities, and the words of $(N_q,b+\kappa_q)$ and of
$\mathrm{BN}_q$ agree.
\end{proof}

\begin{remark}[The corner term]\label{rem:corner-term}
The arrow that $\kappa_q$ deletes is a chord of the closed component of $(N_q,b)$ around the corner
$(0,0)$.  Deleting it opens that component into an arc with two ends at $(0,0)$, and the added arrow
joins one of them to the end of $\alpha_q$.  Thus $\kappa_q$ changes which two of the three arc ends
at $(0,0)$ are joined.  An exhaustive search over the complexes obtained from
$(N_q,b)$ by deleting one arrow and adding one arrow between two generators that meet at most one
arrow finds, for each of $b_{S_{18}}$, $b_{S_{25}}$ and $b_5$, exactly two that are strictly
isomorphic to $\mathrm{BN}_q$ (566 candidates for each $b$ at $q=7$, and 402 at $q=5$).  In each, the
added arrow is labelled $R$ and starts at the end of $N_q$ at $(0,0)$; $\kappa_7$ is the one common
to $b_{S_{18}}$ and $b_{S_{25}}$, and the others are listed in \S\ref{ss:corner-terms}.

Call an element of $\operatorname{End}(X)$ an \emph{algebraic switch} of a reduced twisted complex
$X$ if it replaces two arrows $i\xrightarrow{Y^k}j$ and $i'\xrightarrow{Y^k}j'$ of $\delta_X$, with
$j$ and $j'$ of the same idempotent, by $i\xrightarrow{Y^k}j'$ and $i'\xrightarrow{Y^k}j$, and if the
result satisfies the Maurer--Cartan equation.  Of the $45$ distinct switches of census smoothings at
$q=7$ (\S\ref{ss:smoothings}), $41$ are algebraic switches; each of the other four exchanges the heads of
an arrow labelled $M$ and an arrow labelled $M^2$.  The term $\kappa_q$ is not a sum of switches: no sum
of at most three algebraic switches of $N_q$ is strictly isomorphic to $\mathrm{BN}_q$ (there are
$110$ algebraic switches of $N_7$ and $221\,925$ such sums, and $58$ switches of $N_5$ and $32\,567$
sums), and no sum of at most three of the $45$ census switches is strictly isomorphic to $\mathrm{BN}_7$.
\end{remark}

\begin{remark}[The opposite sign of the perturbation]\label{rem:sign-positive}
By Remark~\ref{rem:sign-pairings}, the pairings for small $t>0$ are those of the present objects with
the curves $\Phi(E_r)$.  Since $\Phi$ is an involution, $\Phi(E_r)=\widetilde{\Kh}(Q_{-r})$, and the
comparison object becomes $\Phi(\mathrm{BN}_q)=\widetilde{\mathrm{BN}}(T_q)$.  The curves $\alpha_q$
and $c$ are $\Phi$-invariant, so part (a) is unchanged; in (b) the slide takes place at the end
$(\pi,0)$; and (c) holds with
\[
 \kappa'_7=2\xrightarrow M3+18\xrightarrow{M^2}3,\qquad \kappa'_5=10\xrightarrow M11+22\xrightarrow{M^2}11
\]
in place of $\kappa_7$ and $\kappa_5$, where $18$ and $22$ are the ends at $(\pi,0)$.  The two corner
terms are not interchangeable.  For $t<0$, the object $(N_7,b_{S_{18}}+\kappa'_7)$ pairs to $13$
with $E_{-1}$, whereas $\rk\widetilde{\Kh}(K_{-1}(7))=11$; under the CHKK correspondence this
pairing would compute $\dim\Inat(K_{-1}(7))$, contradicting $\dim\Inat\le\rk\widetilde{\Kh}$
\cite{KM2}.  At $q=5$ the corresponding values are $9$ and $7$.
\end{remark}

\subsection{Pairing with linear curves}\label{ss:kh-pairing}
For curves $X$, $Y$ in $\Pst$, at least one of them compact, let $i(X,Y)$ be the minimal number of
intersection points of curves homotopic to $X$ and $Y$, where arcs are homotoped through arcs with
the same ends.

\begin{definition}\label{def:linear}
Let $w\in\Z^2$ be primitive.  A compact curve $R\subset\Pst$ is \emph{linear of direction $w$} if every
lift of $R$ lies within a small distance of a line of direction $w$ and is a graph over it, and
either the line meets no lattice point, in which case $R$ is the simple closed curve of direction
$w$, or the line is a lattice line and the lift passes each of its lattice points on a prescribed
side, the choices of side being invariant under $G$.
\end{definition}

Every simple closed curve is linear.  So is every earring $E_r$, whose lifts pass the consecutive
lattice points of the lattice line containing the lifts of its arc on alternating sides, and more
generally every figure-eight around a straight arc between two corners.

\begin{theorem}[Theorem~\ref{thm:pairing}]\label{thm:pairing-full}
Let $5\le q\le21$ be odd, and let $R$ be a linear curve whose direction is parallel neither to
$\alpha_q$ nor to $c$.
\begin{enumerate}
\item[(a)] $i(R,\mathrm{BN}_q)=i(R,\alpha_q)+2\,i(R,c)$.
\item[(b)] If $R$ is bigraded, then
 \[
  \dim\HF(R,\mathrm{BN}_q)=\dim\HF(R,\alpha_q)+2\dim\HF(R,c)=\dim\HF\bigl(R,\alpha_q\oplus(c,J_2)\bigr).
 \]
\end{enumerate}
\end{theorem}

\begin{proof}
\emph{Part (b) from part (a).}  For two $\delta$-graded compact curves that are not parallel, with
local systems $X$ and $X'$, the dimension of $\HF$ is $(\dim X)(\dim X')$ times their minimal
intersection number \cite[Thm.~4.45]{Zibrowius}; this applies to $R$ and $(c,J_2)$ and gives
$\dim\HF(R,(c,J_2))=2\,i(R,c)=2\dim\HF(R,c)$.  For bigraded curves one of which is an arc, the same
formula is \cite[Thm.~5.25]{KWZ}, applied after moving all curves back by $\Phi$; Kotelskiy, Watson
and Zibrowius write out the proof for compact curves, which is the proof of
\cite[Thm.~4.45]{Zibrowius}, and state that the remaining cases are treated similarly.  It applies to
$R$ and the bigraded arcs $\alpha_q$ and $\mathrm{BN}_q$.  Hence (a) implies (b).

\emph{Geodesic representatives.}  Fix a small $\varepsilon>0$ and remove from $\R^2$ the open
$\varepsilon$-discs $B(p)$ about the lattice points $p$.  Their boundary circles are locally
concave, so the flat surface that remains is locally CAT(0) by the theorem of Alexander, Berg and
Bishop \cite[Thm.~7.28]{AKP}, and its universal cover is CAT(0) by the globalization theorem
\cite[Thm.~5.6]{AKP}.  Two geodesics in a CAT(0) space meet in a connected set.  Hence geodesic
representatives in the quotient by $G$, in which an arc ends on the boundary circle of its corner,
realize the minimal intersection numbers, provided a common arc of tangency along a boundary circle
is counted once if the four ends of the two curves near it interleave and not at all otherwise.  The
representatives are: for $R$, the lines of direction $w$ tangent to the discs of their lattice points
on the prescribed sides; for $\alpha_q$, the part of $[p_M,p_0]$ outside the two discs; for $c$, the
line in $\ell^+$ parallel to $\ell$ at distance $\varepsilon$; and for $\mathrm{BN}_q$, the segment
from $\partial B(p_M)$ to a tangent point of $\partial B(p_0)$, an arc of $\partial B(p_0)$, and then
the same parallel line up to $\partial B(p_4)$, where it ends.  The segment $[p_M,p_0]$ meets no other
disc because its direction is primitive.

\emph{Reduction to one lattice point.}  A lift of an arc has trivial stabilizer, so for an arc $X$
the number $i(R,X)$ counts the intersections of the full preimage of $R$ with one lift of $X$.  A
fundamental domain for the stabilizer of the lift of $c$ is one period, so $2\,i(R,c)$ counts the
intersections with the half-open stretch of the parallel line from $p_0$ to $p_4$, the end near $p_0$
included.  Compare the lift of $\mathrm{BN}_q$ with the union of the lift of $\alpha_q$ and this
stretch.  Outside small neighborhoods of $p_M,p_0,\dots,p_4$ the two coincide.  Near $p_M$ both end
along the same segment, and near $p_1,p_2,p_3$ both pass straight.  The translation by
$p_4-p_0\in2\pi\Z^2$ belongs to $G$, so it carries the lifts of $R$ through $p_4$ bijectively onto
those through $p_0$, preserving the sides.  It therefore suffices to show, for every lift $\lambda$ of
$R$ through $p_0$, that
\begin{equation}\label{eq:local-identity}
 \operatorname{corner}(a,b')+\operatorname{end}(-b')=\operatorname{end}(a)+1,
\end{equation}
where $a$ and $b'$ are the directions from $p_0$ to $p_M$ and to $p_1$; $\operatorname{end}(x)$ is the
number of intersections of $\lambda$ with an arc that ends on $\partial B(p_0)$ along the direction
$x$; $\operatorname{corner}(a,b')$ is the number of intersections of $\lambda$ with the part of
$\mathrm{BN}_q$ near $p_0$; and the summand $1$ counts the straight passage of the stretch of $c$.

\emph{Local counts.}  Let $D$ be the open half-plane bounded by the lattice line of $\lambda$ into
which $\lambda$ is displaced near $p_0$, and $D'$ the other one.  The directions $a$ and $b'$ are not
parallel to $w$, so each points into $D$ or into $D'$.  The line $\lambda$ is tangent to $\partial
B(p_0)$ on the side of $D$.  An arc ending orthogonally on $\partial B(p_0)$ along $x$ meets it once
if $x$ points into $D$ and not at all otherwise.  A straight passage meets it once, since two distinct
tangent lines of a circle meet in one point.  The angle between $a$ and $b'$ is less than $\pi$, and
the part of $\mathrm{BN}_q$ near $p_0$ runs in along $a$, around the far side of $B(p_0)$ from that
angle, and out along $b'$; it meets $\lambda$ once if $a$ and $b'$ point into different half-planes,
twice if both point into $D$, and not at all if both point into $D'$.  Since
$\operatorname{end}(-b')=1$ exactly when $b'$ points into $D'$, the two sides of
\eqref{eq:local-identity} are
\begin{center}
\begin{tabular}{lcccc}
\toprule
 & $a\in D$, $b'\in D'$ & $a\in D'$, $b'\in D$ & $a,b'\in D$ & $a,b'\in D'$\\
\midrule
left side & $1+1$ & $1+0$ & $2+0$ & $0+1$\\
right side & $1+1$ & $0+1$ & $1+1$ & $0+1$\\
\bottomrule
\end{tabular}
\end{center}
and they agree in every case.  This proves (a).
\end{proof}

\begin{corollary}[Corollary~\ref{cor:closures}]\label{cor:closures-full}
Let $q\in\{5,7\}$ and let $b$ be as in Theorem~\ref{thm:structure-full}(a).  For every
$r\in\Q\cup\{\infty\}$,
\[
 \dim\HF\bigl(E_r,(N_q,b)\bigr)=\dim\HF\bigl(E_r,(N_q,b+\kappa_q)\bigr)=\rk\widetilde{\Kh}\bigl(K_r(q);\F_2\bigr).
\]
\end{corollary}

\begin{proof}
The second equality is Theorem~\ref{thm:structure-full}(c) together with \eqref{eq:kwz-pairing}.  The
earring $E_r$ is linear, and after moving all curves back by $\Phi$ it becomes the bigraded curve
$\widetilde{\Kh}(Q_{-r})$.  If its direction is parallel neither to $\alpha_q$ nor to $c$,
Theorems~\ref{thm:structure-full}(a) and~\ref{thm:pairing-full}(b) give
$\dim\HF(E_r,(N_q,b))=\dim\HF(E_r,\mathrm{BN}_q)$.  The two remaining earrings are $E_{-1}$, parallel
to $c$, and $E_{r_q}$, parallel to $\alpha_q$, and for them all three numbers are computed directly:
\begin{center}
\begin{tabular}{ccccccc}
\toprule
 & \multicolumn{3}{c}{$E_{-1}$} & \multicolumn{3}{c}{$E_{r_q}$}\\
\cmidrule(lr){2-4}\cmidrule(lr){5-7}
$q$ & $(N_q,b)$ & $(N_q,b+\kappa_q)$ & $\rk\widetilde{\Kh}$ & $(N_q,b)$ & $(N_q,b+\kappa_q)$ &
 $\rk\widetilde{\Kh}$\\
\midrule
$5$ & $7$ & $7$ & $7$ & $30$ & $30$ & $30$\\
$7$ & $11$ & $11$ & $11$ & $46$ & $46$ & $46$\\
\bottomrule
\end{tabular}
\end{center}
The undeformed complex $N_q$ pairs to $7$ and $28$ at $q=5$, and to $11$ and $44$ at $q=7$.
\end{proof}

We have also computed the three numbers directly at 340 slopes for $q=5$ and 342 slopes for $q=7$:
every slope with $|r|\le2$ and denominator at most $14$, every slope with $|r|\le5$ and denominator
at most $6$, the integers with $|r|\le10$, $\pm r_q$ and $\infty$.  These include 108 slopes with odd
numerator and odd denominator, for which the arc of $E_r$ joins $(\pi,0)$ to $(0,\pi)$.  The three numbers agree
at every slope.  The undeformed complex $N_q$ pairs to less than $\rk\widetilde{\Kh}$ at 40 of these
slopes for $q=5$ and at 43 for $q=7$, and never to more.  Under the CHKK correspondence the pairing of
the assigned object with $E_r$ computes $\dim\Inat(K_r(q))$, which is at most
$\rk\widetilde{\Kh}(K_r(q))$ \cite{KM2}.  Both $(N_q,b)$ and $(N_q,b+\kappa_q)$ attain this bound at
every rational closure: at $r=-1/2$ the common value is $q+2=\dim\Inat(P(-2,3,q))$
(Theorem~\ref{thm:inat}), and at $r=-3/4$ and $q=7$ it is $31=\dim\Inat(K_*)$
(Proposition~\ref{prop:aux-closure}).

\begin{remark}[The direction of $c$]\label{rem:exceptional}
The hypothesis that the direction is not parallel to $c$ cannot be dropped.  Among figure-eights around
straight arcs and simple closed curves, exactly two curves pair differently with $\mathrm{BN}_q$ and
with $\alpha_q\oplus(c,J_2)$, and both are parallel to $c$: the curve $c$ itself, with $11$ against
$13$ at $q=7$ and $7$ against $9$ at $q=5$, and the figure-eight around the arc parallel to $c$ that
joins $(0,0)$ to $(\pi,\pi)$, with $13$ against $11$ and $9$ against $7$.  The latter surrounds the corner
$(\pi,\pi)$, which no tangle character variety meets, and is not the earring of a rational tangle;
the earring $E_{-1}$, around the arc joining $(\pi,0)$ to $(0,\pi)$, pairs equally.  For general side
patterns we have computed all linear curves with direction $(m,n)$, $|m|,|n|\le6$, or parallel to
$\alpha_q$, on lattice lines of both orbits and with every side pattern of period $2$, $4$ or $6$,
together with the simple closed curves: $1887$ curves for each $q=5,7,\dots,13$.  For each $q$,
exactly $22$ of them pair differently, all parallel to $c$: all $19$ side patterns on the lattice line
through $(0,0)$ and $(\pi,\pi)$, the two constant patterns on the line through $(\pi,0)$ and
$(0,\pi)$, and the simple closed curve.  The constant patterns on either line and the simple closed
curve are isotopic to $c$.  No curve parallel to
$\alpha_q$ pairs differently, although the proof does not cover that direction.
\end{remark}

\begin{remark}[Curves parallel to $c$]\label{rem:sharpness}
Curves parallel to $c$ separate $(N_q,b)$ from $(N_q,b+\kappa_q)\cong\mathrm{BN}_q$.  At $q=7$ the
curve $c$ pairs to $11$ with $\mathrm{BN}_7$ and to $13$ with $(N_7,b_{S_{18}})$.  The curve $(c,J_2)$
pairs to $22$ and to $26$, the difference being $\dim\HF\bigl((c,J_2),(c,J_2)\bigr)=4$.  The image
under $\Phi$ of the special component of $\widetilde{\Kh}(T_7)$ in the sense of \cite{KWZ}, the
component other than the figure-eight $\widetilde{\Kh}(Q_{-r_7})$, is parallel to $c$ and pairs to
$52$ and to $44$.  At $q=5$ the corresponding values are $7$, $9$; $14$, $18$; and $36$, $28$.
\end{remark}

\subsection{The instanton object of $T_q$}\label{ss:kh-question}
Theorem~\ref{thm:structure-full} and Corollary~\ref{cor:closures-full} suggest, without deciding, the
following question.

\begin{question}[Question~\ref{q:arc}]\label{q:arc-full}
Let $q\ge5$ be odd with $\gcd(q,3)=1$.  Suppose that the CHKK tangle assignment is defined for $T_q$
with Smith's perturbation for small $t<0$, and let $X_q$ be the object of $\operatorname{Tw}(\mathcal
B)$ that it determines, an immersed curve with a bounding cochain, drawn in Smith's coordinates.  Is
$X_q$ homotopy equivalent to $\mathrm{BN}_q$?
\end{question}

For $q=5$ and $q=7$ the evidence is as follows.  The object $\mathrm{BN}_q$ arises from the
piecewise-linear model of Smith's curve by the explicit Maurer--Cartan element $b+\kappa_q$
(Theorem~\ref{thm:structure-full}).  By Corollary~\ref{cor:closures-full}, $(N_q,b)$ and
$(N_q,b+\kappa_q)$ have the same pairing with every rational earring, equal to the rank of reduced
Khovanov homology of the closure; both are compatible with $\dim\Inat\le\rk\widetilde{\Kh}$ at every
rational closure, and rational closures cannot distinguish them.  They are distinguished by curves
parallel to $c$ (Remark~\ref{rem:sharpness}), and by gradings: $\mathrm{BN}_q$ is bigraded
and $(N_q,b)$ is not, but no quantum grading is known on the instanton side, so this difference is not
a test.

\begin{proposition}\label{prop:hyp-conflict}
The object $\mathrm{BN}_7$ is not homotopy equivalent to any object allowed by
Hypothesis~\ref{hyp:q7-support}, although $\dim\HF(E_r,\mathrm{BN}_7)=\dim\HF(E_r,(N,b_{S_{18}}))$ for
every $r\in\Q\cup\{\infty\}$.
\end{proposition}

\begin{proof}
The dimension of $\HF(Z,-)$ for a fixed object $Z$ is invariant under homotopy equivalence.  By
\eqref{eq:kwz-pairing}, $\mathrm{BN}_7$ pairs to $9$ with $E_{-1/2}$ and to $31$ with $E_{-3/4}$, the
ranks of reduced Khovanov homology of $P(-2,3,7)$ and $K_*$ (Section~\ref{sec:inat} and
Proposition~\ref{prop:aux-closure}).
Among the objects allowed by Hypothesis~\ref{hyp:q7-support}, only $(N,b_{S_{18}})$ and
$(N,b_{S_{25}})$ have this pair (proof of Corollary~\ref{cor:aux-conditional}), and they are strictly
isomorphic (Corollary~\ref{cor:q7-classes}).  By Theorem~\ref{thm:structure-full}(a),
\[
 \dim\HF\bigl((c,J_2),(N,b_{S_{18}})\bigr)=\dim\HF\bigl((c,J_2),\alpha_7\bigr)
 +\dim\HF\bigl((c,J_2),(c,J_2)\bigr)=22+4=26,
\]
whereas $\dim\HF\bigl((c,J_2),\mathrm{BN}_7\bigr)=22$.  The last statement is
Corollary~\ref{cor:closures-full}.
\end{proof}

Consequently an affirmative answer to Question~\ref{q:arc-full} for $q=7$ would show that
Hypothesis~\ref{hyp:q7-support} fails, and no rational closure could detect the failure: the two
closures of Corollary~\ref{cor:aux-conditional}, like every other rational closure, assign the same
dimension to $\mathrm{BN}_7$ as to the class of $(N,b_{S_{18}})$ that the hypothesis selects.

\begin{remark}[The cases $q\ge11$]\label{rem:q11}
For $q=11$ and $13$ the encoding of a numerical trace of Smith's curve is strictly isomorphic to
$N_q$, with $47$ and $55$ generators, and its pairing with $E_{-1/2}$ is $q+4$, two more than
$\rk\widetilde{\Kh}(P(-2,3,q))=q+2$; for $q=17$ and $19$ the pairing of the traced curve is $q+8$.  At
$q=5,7$ the pairing is $q$, and one switch raises it to $q+2$.  The deformation of
Theorem~\ref{thm:structure-full} does not persist in the following finite sense.  Call the minimal
number of arrows that must be deleted from $\delta_{N_q}$, and the same number added, to reach a
complex strictly isomorphic to $\mathrm{BN}_q$ the \emph{edit distance}.  It is $2$ for $q=5,7$, $4$
for $q=11,13$, and at least $6$ for $q=17,19$.  At $q=5,7$ one algebraic switch followed by one
replacement of an arrow reaches $\mathrm{BN}_q$; at $q=11$ and $13$ no chain of at most three
algebraic switches, each Maurer--Cartan for the complex it deforms, ends one replacement away from
$\mathrm{BN}_q$.  At $q=11$ the first such chains have four switches, and they end at the sum of an
arc and a closed curve other than $\alpha_{11}\oplus(c,J_2)$; the latter is more than five switches
away from $N_{11}$, against one at $q=5,7$.  None of the $35\,004$ complexes ($6\,331$ up to strict
isomorphism) obtained from $N_{11}$ by at most two algebraic switches has $\alpha_{11}$ as a summand, and each of them either
pairs with $E_{-1/2}$ to a number other than $\dim\Inat(P(-2,3,11))=13$, or pairs with some earring
of a panel of $25$ slopes to more than the rank of reduced Khovanov homology of the closure.  These
are finite searches within the class of algebraic switches, and they do not exclude a larger bounding
cochain.
\end{remark}
